% =====================================================================
\section{Results}
\label{sec:results}
% Sources (all committed, all regenerated without model calls):
%   core            backend/out_study/T2/analysis.json (identity to the clean
%                   output and distinct outputs per cell: its "exploratory"
%                   block), T2/o5_sensitivity.json
%   C1              T2C1/analysis.json        placement  T2P/analysis.json
%   severity        T2S/analysis.json         warm       T2W/analysis.json
%   real defects    T1/analysis.json          downstream T2D/analysis.json
%   overwrite       T2D19/analysis.json
%   second model    T2L/analysis.json (repair_analysis.py --cell c7),
%                   C7_preflight_r3/gates.json, T2L/verification.json
%   origin, class ordering, change->resolve, calibration:
%                   backend/out_study/secondary_extra.json
%   summary record  docs/paper/RESULTS.md
% Rates are per-defect means of the r = 3 repetitions; intervals are 95%
% bootstrap intervals over defects unless a line says otherwise.
% =====================================================================

Unless stated otherwise, rates are means over defects of the per-defect
average across repetitions, and intervals are 95\% bootstrap intervals that
resample defects. Differences are in percentage points (pp).

% ---------------------------------------------------------------------
\subsection{RQ1: Feedback content does not change repair}
\label{sec:results-rq1}

The 18 excluded trials cost three defects their \arm{A1}/\arm{A2} pair. The
remaining $N=139$ defects completed all three repetitions in both arms and form
the primary contrast. Measured-fact feedback repaired \textbf{94.7\%}
of defects and the contentless rejection \textbf{94.0\%}. The difference is
\textbf{+0.7\pp{}, 95\% CI \ci{-2.9}{+4.3}}. On the majority dichotomy, \arm{A1}
repaired 131 defects and \arm{A2} 129. The discordant pairs split 6 to 4, and
the exact McNemar test gives $p=0.754$. The confirmatory test does not reject.
The interval rules out an advantage for measured facts larger than 4.3\pp{}, and
a disadvantage larger than 2.9\pp{}, on this model and these defects. The
publication rule of Section~\ref{sec:design-analysis} required at least 12\pp{},
so the repair study did not become a paper of its own.

\begin{table*}[t]
\centering
\caption{The confirmatory contrast and its pre-specified companions (core
cell). Rates are per-defect means over $r=3$. $p$ is the exact McNemar test on
the majority dichotomy. Only the first row is confirmatory.}
\label{tab:rq1}
\footnotesize
\begin{tabularx}{\linewidth}{@{}Y r r r l r@{}}
\toprule
Contrast & $N$ & Rate 1 & Rate 2 & Difference [95\% CI] & $p$ \\
\midrule
\textbf{\arm{A1}\,$-$\,\arm{A2}, primary} & \textbf{139} & \textbf{94.7} &
  \textbf{94.0} & \textbf{+0.7 \ci{-2.9}{+4.3}} & \textbf{0.754} \\
\quad without multi-call trials & 137 & 95.4 & 94.4 & +1.0 \ci{-2.2}{+4.1} &
  0.727 \\
\quad four \op{O5} matchers$^a$ & 139 & 94.7 & 94.0 & +0.7
  \ci{-2.9}{+4.3} & 0.754 \\
\quad without \op{O5} & 120 & 95.0 & 95.0 & 0.0 \ci{-3.6}{+3.6} & 1.000 \\
\quad \op{O1} only (production-like) & 24 & 84.7 & 84.7 & 0.0
  \ci{-12.5}{+12.5} & 1.000 \\
\quad severe cell only (\op{O2}--\op{O7}) & 115 & 96.8 & 95.9 & +0.9
  \ci{-2.6}{+4.6} & 0.727 \\
\midrule
\arm{A1}/\arm{A0} (ratio) & 71 & & & 0.96 \ci{0.92}{1.00} & -- \\
\arm{A2}\,$-$\,\arm{A3} & 34 & 96.1 & 92.2 & +3.9 \ci{-4.9}{+13.7} & 1.000 \\
Cochran's $Q$, \arm{A1}/\arm{A2}/\arm{A3} & 34 & \multicolumn{3}{l}{$Q=0.50$,
  df\,$=2$} & 0.779 \\
\bottomrule
\end{tabularx}

\smallskip
\parbox{\linewidth}{\footnotesize $^a$ Content-word Jaccard $\ge 0.10$
(primary), $\ge 0.05$, containment $>0$, and the production character-level
matcher. All four give the same contrast to the digit, and recomputing the
primary reproduces all 157 stored \op{O5} outcomes.}
\end{table*}

Table~\ref{tab:rq1} shows that the null does not depend on any analytic
choice. It holds without the 47 trials that needed a schema-correction second
call. It holds under every \op{O5} matcher and without \op{O5} altogether. It
holds within the production-like \op{O1} cell, where the interval is wide
because $N=24$, and within the severe cell. Measured-fact feedback recovers
96\% of the capability ceiling (\arm{A1}/\arm{A0} = 0.96 \ci{0.92}{1.00}).
\arm{A0} repaired every one of its 209 trials.
Figure~\ref{fig:forest} sets the primary contrast beside selected
\arm{A1}\,$-$\,\arm{A2} contrast the study measured, including the
exploratory exact-restoration contrast and the second-model replication.
On overall repair, every interval measured on the hosted model includes zero.
Only the second model's excludes it (Section~\ref{sec:results-rq6}), alongside
the exploratory exact-restoration outcome (Section~\ref{sec:results-rq2}).

Repair rarely went wrong in any other way. Of 421 usable \arm{A1} trials, 399
were \cls{repaired}, 16 \cls{replaced}, 3 \cls{deleted} and 3
\cls{unrepaired}. \arm{A2}'s 417 trials split 392, 20, 3 and 2. The supervisor
accepted 415 of the 421 retried \arm{A1} outputs (\verdict{proceed} or
\verdict{proceed\_with\_warning}) and 410 of the 417 \arm{A2} outputs.
Collateral damage was 1.5\% of measurable \arm{A1} trials and 0.4\% for
\arm{A2}. For \arm{A3}, which regenerates from scratch, it was 16.7\%.

\begin{figure*}[t]
\centering
% =====================================================================
%  Fig. forest -- selected A1 - A2 contrasts (pp), with 95%
%  bootstrap intervals over defects. Values transcribed from:
%    rows 1-5   backend/out_study/T2/analysis.json  (confirmatory + secondary)
%    rows 6-7   backend/out_study/T2C1/analysis.json (gap_C0, gap_C1)
%    rows 8-9   backend/out_study/T2W/analysis.json  (A1@T0-A2@T0, A1@T1-A2@T1)
%    row 10     backend/out_study/T2D/analysis.json  (E1_contrasts.A1-A2)
%    row 11     backend/out_study/T2L/analysis.json  (A1_minus_A2; second
%               model, amendments A10/A11 -- the only non-hosted row)
%  Colours: dataviz reference palette -- slot 1 (#2A78D6) for the one
%  confirmatory row, secondary ink (#52514E) for the rest (emphasis scheme);
%  solid hairline grid; text stays in ink colours, never the mark colour.
%  Width budget: labels ~4.3 cm + plot 5.6 cm + estimates ~2.9 cm < 18.4 cm
%  (the two-column layout's full-width figure* text block).
% =====================================================================
\begin{tikzpicture}[x=0.135cm, y=-0.5cm, font=\footnotesize]
  \definecolor{fpaccent}{HTML}{2A78D6}
  \definecolor{fpink}{HTML}{52514E}
  \definecolor{fpgrid}{HTML}{E4E3DE}
  \definecolor{fpzero}{HTML}{8F8E89}

  % \fprow{y}{label}{lo}{est}{hi}{text}{colour}
  \newcommand{\fprow}[7]{%
    \node[anchor=east, align=right] at (-15.5,#1) {#2};
    \draw[#7, line width=0.9pt] (#3,#1) -- (#5,#1);
    \filldraw[fill=#7, draw=white, line width=0.7pt] (#4,#1) circle (2.4pt);
    \node[anchor=west, text=fpink, font=\scriptsize] at (22.5,#1) {#6};
  }

  % last data row (exact restoration) sits at y = 15.0
  \def\ybottom{15.8}

  % primary interval band
  \fill[fpaccent!10] (-2.9,0.35) rectangle (4.3,\ybottom);
  % gridlines and zero line (solid hairlines)
  \foreach \gx in {-10,-5,5,10,15,20}
    \draw[fpgrid, line width=0.35pt] (\gx,0.35) -- (\gx,\ybottom);
  \draw[fpzero, line width=0.5pt] (0,0.35) -- (0,\ybottom);

  % column heads
  \node[anchor=east, font=\footnotesize\bfseries] at (-15.5,-0.4)
       {Contrast ($N$ defects)};
  \node[anchor=west, font=\scriptsize\bfseries] at (22.5,-0.4)
       {Estimate [95\% CI]};

  % --- core (confirmatory, then pre-specified companions)
  \fprow{1}{\textbf{Core, primary} (139)}{-2.9}{0.7}{4.3}%
        {\textcolor{black}{\textbf{+0.7 [$-$2.9, +4.3]}}}{fpaccent}
  \fprow{2}{Core, no multi-call trials (137)}{-2.2}{1.0}{4.1}%
        {+1.0 [$-$2.2, +4.1]}{fpink}
  \fprow{3}{Core, without \op{O5} (120)}{-3.6}{0.0}{3.6}%
        {\phantom{+}0.0 [$-$3.6, +3.6]}{fpink}
  \fprow{4}{Core, \op{O1} only (24)}{-12.5}{0.0}{12.5}%
        {\phantom{+}0.0 [$-$12.5, +12.5]}{fpink}
  \fprow{5}{Core, severe cell (115)}{-2.6}{0.9}{4.6}%
        {+0.9 [$-$2.6, +4.6]}{fpink}
  % --- constraint disclosure
  \fprow{6.6}{C1 cell, production prompt (37)}{-6.3}{-1.8}{1.8}%
        {$-$1.8 [$-$6.3, +1.8]}{fpink}
  \fprow{7.6}{C1 cell, obligations stated (37)}{-6.3}{-1.8}{1.8}%
        {$-$1.8 [$-$6.3, +1.8]}{fpink}
  % --- warm resampling
  \fprow{9.2}{Warm cell, $T{=}0$ (45)}{-7.4}{-3.0}{0.0}%
        {$-$3.0 [$-$7.4, \phantom{+}0.0]}{fpink}
  \fprow{10.2}{Warm cell, $T{=}1$ (45)}{0.0}{2.2}{5.2}%
        {+2.2 [\phantom{$-$}0.0, +5.2]}{fpink}
  % --- pipeline level
  \fprow{11.8}{Downstream acceptance (88)}{-5.7}{0.0}{5.7}%
        {\phantom{+}0.0 [$-$5.7, +5.7]}{fpink}
  % --- second model (Qwen2.5-7B-Instruct, local)
  \fprow{13.4}{Second model (138)}{0.5}{6.8}{13.3}%
        {+6.8 [+0.5, +13.3]}{fpink}
  % --- exact restoration (exploratory mechanism)
  \fprow{15.0}{\textbf{Exact restoration, core} (139)}{8.2}{14.6}{21.3}%
        {\textcolor{black}{\textbf{+14.6 [+8.2, +21.3]}}}{fpaccent}

  % x axis
  \draw[fpzero, line width=0.4pt] (-15,\ybottom) -- (22,\ybottom);
  \foreach \gx/\lab in {-10/$-$10,-5/$-$5,0/0,5/+5,10/+10,15/+15,20/+20} {
    \draw[fpzero, line width=0.4pt] (\gx,\ybottom) -- (\gx,\ybottom+0.18);
    \node[anchor=north, text=fpink] at (\gx,\ybottom+0.2) {\lab};
  }
  \node[anchor=north] at (5,\ybottom+0.95)
       {\arm{A1}\,$-$\,\arm{A2} (percentage points)};
  \node[anchor=north east, text=fpink, font=\scriptsize] at (4.4,\ybottom+1.75)
       {$\leftarrow$ rejection better};
  \node[anchor=north west, text=fpink, font=\scriptsize] at (5.6,\ybottom+1.75)
       {measured facts better $\rightarrow$};
\end{tikzpicture}

\caption{Selected \arm{A1}\,$-$\,\arm{A2} contrasts, with 95\%
bootstrap intervals over defects. The shaded band is the primary interval. Rows
below the first are exploratory or secondary. The penultimate row is the second
model (Section~\ref{sec:results-rq6}); the bottom row is the exploratory
exact-restoration contrast (Section~\ref{sec:results-rq2}).}
\label{fig:forest}
\end{figure*}

% ---------------------------------------------------------------------
\subsection{Pre-recorded predictions and input size}
\label{sec:results-predictions}
Repair exceeded the pre-recorded predictions for both feedback arms, and the
predicted class ordering reversed: referential defects were hardest. An
exploratory review-count analysis found lower repair or exact restoration at
the largest input stratum for some reasoning operators. The small strata and
nonmonotonic patterns do not establish an input-size mechanism; detailed
predictions, operator-specific rates and slopes are in Appendix~\ref{app:diagnostics}.

% ---------------------------------------------------------------------
\subsection{RQ2: Exact restoration and artifact presentation}
\label{sec:results-rq2}

Many repaired outputs restore the clean output. We define \emph{exact restoration} once for the entire paper as
canonical-JSON identity to the clean pre-injection stage output (parsed JSON
serialised with sorted keys and invariant whitespace). Under this definition,
100\% of trials that achieve exact restoration are verified as \cls{repaired}
(706 of 706 across all \cell{T2} arms, 337 of 337 in \cell{T2L}, and 319 of
319 in \cell{T2W}).

\paragraph{Unconditional exact restoration contrast}
Taking exact restoration as an unconditional outcome over all trials yields a
causally clean arm contrast within matched defects. On the core (\cell{T2},
$N=139$ paired defects), measured facts substantially increase exact
restoration: \arm{A1} restores 68.1\% of defects and \arm{A2} 53.5\%, an
advantage of \textbf{+14.6\pp{} \ci{+8.2}{+21.3}} (exact McNemar $p=0.00132$,
26 \arm{A1}-only vs 7 \arm{A2}-only discordant defects). At the trial level,
288 of 421 \arm{A1} retries (68.4\%) reproduce the clean output as canonical JSON,
compared to 223 of 417 (53.5\%) for \arm{A2}, 149 of 209 (71.3\%) for \arm{A0},
and 46 of 102 (45.1\%) for \arm{A3}.

The exploratory restoration contrast was positive in both warm-cell temperatures; effort-level results vary and are reported in Section~\ref{sec:results-thinking}.
In the warm-resampling cell (\cell{T2W}, $N=45$), \arm{A1} leads \arm{A2} by
\textbf{+20.7\pp{} \ci{+8.1}{+34.1}} at $T=0$ (69.6\% vs 48.9\%, $p=0.021$)
and by \textbf{+20.0\pp{} \ci{+8.9}{+32.6}} at $T=1$ (68.9\% vs 48.9\%,
$p=0.0039$). On the second
model (\cell{T2L}, Qwen2.5-7B, $N=138$), pooled exact restoration is lower and the
overall gap is narrow: \arm{A1} 42.0\% vs \arm{A2} 39.4\% (+2.7\pp{} \ci{-4.6}{+9.9},
$p=0.572$, 16 vs 12 discordant).

\paragraph{Trial-level decomposition}
The decomposition in Appendix~\ref{app:diagnostics} separates exact copies,
repaired rewrites, replacements or deletions, and outputs where the check
still fires. On the local model, the latter occurred in 32/414 \arm{A1} and
107/413 \arm{A2} trials. These outcome categories do not reveal whether the
model located the defect. Within-rewrite comparisons condition on a
post-treatment outcome and are descriptive.

\paragraph{Evidence consistent with copy-and-correct}
At temperature~0, canonical-JSON identity alone cannot separate two mechanisms. The model
might copy the artifact it is shown and fix the one injected change, or it
might re-derive its original answer from the same prompt. The warm-resampling
cell provides evidence bearing on them (Table~\ref{tab:mechanism}). At \arm{P1}, where the
artifact is shown, canonical-JSON identity to the clean output does not move with
temperature: \arm{A1} 69.6\% at $T=0$ and 68.9\% at $T=1$; \arm{A2} 48.9\% at
both. \arm{A3}, which sees no artifact, falls from 45\% identical as canonical JSON in the
core to \textbf{0\%} at $T=1$. A fresh draw at $T=1$ never reproduces the clean
output. An \arm{A1} retry shown the defective output reproduces it about two
times in three, at either temperature. This pattern is consistent with copying and correcting the shown artifact, although it does not identify the model's internal process. Repair stays at the ceiling at $T=1$
(\arm{A1} 99.3\%, \arm{A2} 97.0\%), so the null is not an artefact of greedy
decoding. Neither reading fixed in advance for this cell applies: canonical-JSON identity
did not fall, and neither did repair.

The cell also gives a budget-accounted comparison that
\citet{olausson2024selfrepair} require. \arm{A1}'s single retry at $T=0$ repairs
97.0\% at about 9.7k tokens per trial. \arm{A3}'s best-of-$k$ at $T=1$, taking
repetitions in execution order, reaches 82.2\%, 88.9\% and 93.3\% for $k=1,2,3$,
at about 13.6k tokens per draw (41k for $k=3$). In this setting, one retry with
the artifact in view repairs at least as often as three fresh draws (97.0\%
against 93.3\%), at under a quarter of their token cost. The
\arm{A1}\,$-$\,\arm{A2} gap is $-3.0$\pp{} \ci{-7.4}{0.0} at $T=0$ and
$+2.2$\pp{} \ci{0.0}{+5.2} at $T=1$. Their difference, $+5.2$\pp{}
\ci{+0.7}{+11.1}, is one of many exploratory intervals in this paper and rests
on a handful of defects. We do not read it as evidence that content matters
under sampling.

Repetitions at $T=0$ are nearly identical. Of the \arm{A1} cells
(defect\,$\times$\,arm) with all three repetitions, 73\% produce one distinct
output and 27\% produce two. The figures are 60\% and 40\% for \arm{A2}, where
one cell of 139 gave three distinct outputs, and 77\% and 23\% for \arm{A0}.
The exception is \arm{A3} at \arm{P0}, where 94\% of cells give two distinct
outputs.

\begin{table*}[t]
\centering
\caption{What moves repair (RQ2). Every row is \arm{A1} at \arm{P1} and
$T=0$ unless stated. Each cell re-runs its comparators in-session. Rates are
per-defect means with 95\% intervals over defects. Canonical-JSON identity is per
trial.}
\label{tab:mechanism}
\footnotesize
\begin{tabularx}{\linewidth}{@{}Y r L{2.7cm} L{3.6cm}@{}}
\toprule
Cell and level & $N$ & Repaired [95\% CI] & Contrast [95\% CI] \\
\midrule
\multicolumn{4}{@{}l}{\emph{Placement}} \\
\quad \arm{P0} artifact absent (production) & 45 & 94.1 \ci{88.1}{98.5} &
  \arm{P1}\,$-$\,\arm{P0} +3.7 \ci{-1.5}{+9.6} \\
\quad \arm{P1} artifact under review & 45 & 97.8 \ci{94.1}{100} &
  \arm{P1}\,$-$\,\arm{P2} 0.0 \ci{-4.4}{+4.4} \\
\quad \arm{P2} artifact as own turn & 45 & 97.8 \ci{93.3}{100} &
  \arm{P3}\,$-$\,\arm{P2} $-0.7$ \ci{-2.2}{0.0} \\
\quad \arm{P3} own turn + marker & 45 & 97.0 \ci{91.1}{100} &
  \arm{P1}\,$-$\,\arm{P3} +0.7 \ci{-2.2}{+4.4} \\
\multicolumn{4}{@{}l}{\emph{Temperature (warm cell)}} \\
\quad \arm{A1}, $T=1$ & 45 & 99.3 \ci{97.8}{100} & identity 68.9\% (69.6\% at
  $T=0$) \\
\quad \arm{A2}, $T=1$ & 45 & 97.0 \ci{94.1}{99.3} & identity 48.9\% (48.9\% at
  $T=0$) \\
\quad \arm{A3} at \arm{P0}, $T=1$ & 45 & 80.7 \ci{71.9}{88.9} & identity 0\%
  (45\% at $T=0$, core) \\
\multicolumn{4}{@{}l}{\emph{\op{O1} severity}} \\
\quad level 1: one id corrupted & 24 & 86.1 \ci{72.2}{97.2} & 1\,$-$\,2:
  $-2.8$ \ci{-11.1}{+5.6} \\
\quad level 2: every id of the pattern & 24 & 88.9 \ci{76.4}{98.6} & \\
\multicolumn{4}{@{}l}{\emph{Constraint disclosure (gap = \arm{A1}\,$-$\,\arm{A2})}} \\
\quad C0, production prompt & 37 & \arm{A1} 94.6, \arm{A2} 96.4 & gap $-1.8$
  \ci{-6.3}{+1.8} \\
\quad C1, obligations restated & 37 & \arm{A1} 92.8, \arm{A2} 94.6 & gap $-1.8$
  \ci{-6.3}{+1.8} \\
\quad difference of gaps & 37 & & 0.0 \ci{-6.3}{+6.3} \\
\bottomrule
\end{tabularx}
\end{table*}

\paragraph{Placement}
With the artifact in view, the observed framing contrasts were small. \arm{P1}, \arm{P2} and
\arm{P3} repair at 97.0--97.8\%, and the paired contrasts among them are within
a defect or two of zero. The pre-recorded prediction
\arm{P1}\,$\gg$\,\arm{P2} does not hold. A self-attribution effect of the kind
reported by \citet{tsui2026scbench} and \citet{chen2026rolerelabel} cannot
appear at a ceiling. \arm{P0}, the production placement, repairs 94.1\%, which
is not detectably lower than \arm{P1} ($+3.7$\pp{} \ci{-1.5}{+9.6}). It differs
in two other ways. First, \arm{P0} retries damage the rest of the output:
collateral damage is \textbf{38.4\%} of measurable trials at \arm{P0}, against
1.1\% at \arm{P1} and 3.4\% at \arm{P2} and \arm{P3}. A model that cannot see
its output regenerates everything, and the regeneration churns items that were
correct. Second, every one of the 16 report-stage defects needed a
schema-correction second call at \arm{P0}, and 50 of the cell's 52 multi-call
trials were at \arm{P0}. Without its multi-call trials, \arm{P0} repairs 90.8\%
\ci{81.6}{97.7} of the remaining 29 defects.

\paragraph{Severity and constraint disclosure}
Breaking all of a pattern's evidence did not lower repair. Level~2 of \op{O1}
repaired 88.9\%, against 86.1\% at level~1 ($-2.8$\pp{} \ci{-11.1}{+5.6}).
Restating the stage's obligations at retry time yielded no detectable difference in this small cell. The
\arm{A1}\,$-$\,\arm{A2} gap was $-1.8$\pp{} under both C0 and C1, and the
difference of gaps was 0.0\pp{} \ci{-6.3}{+6.3}, as pre-recorded. This imprecise comparison does not rule out an effect of disclosure.

\paragraph{Serving drift}
Every exploratory cell re-ran its comparators in-session, so each re-run is a
test--retest measurement of the hosted model at temperature~0, days after the
core. Against the core's trials for the same defect, arm and repetition, the
re-runs reproduced the outcome class in 95.8--99.6\% of pairs, across the four
cells that re-ran a core condition. They reproduced the output byte for byte in
only 76.4--83.3\% of pairs. For the C1 cell, 213 of 222 pairs (96\%) kept their
class and 185 (83\%) were byte-identical. Temperature~0 on a hosted API is
reproducible in outcome, not in output. A comparison across sessions would have
mixed that drift into every contrast. A local runtime was steadier, but not
bit-identical either (Section~\ref{sec:results-rq6}).

% ---------------------------------------------------------------------
\subsection{RQ3: Real defects are not injected defects}
\label{sec:results-rq3}

\begin{table*}[t]
\centering
\caption{The production retry on real, production-caught defects (13 defects in
7 business\,$\times$\,stage cells; arms at \arm{P0}; $r=3$). ``Accepted'' means
the supervisor routes the retried output to no retry. Intervals resample cells.
Descriptive only.}
\label{tab:tier1}
\footnotesize
\begin{tabularx}{\linewidth}{@{}Y r r r@{}}
\toprule
& \arm{A1} (production) & \arm{A2} & \arm{A3} \\
\midrule
Accepted & 28.2 \ci{15.1}{46.2} & 76.9 \ci{55.6}{89.7} & 25.6 \ci{12.8}{41.7}
  \\
Originally fired checks resolved & 38.5 & 76.9 & 33.3 \\
Returned the original defective output & 0.0 & 0.0 & 48.7 \\
Trace class (8 defects): production check accepts & 29.2 & 79.2 & 29.2 \\
Trace class: content-word check accepts & 87.5 & 100.0 & 79.2 \\
\midrule
\multicolumn{4}{@{}l}{\emph{Same criterion, injected defects at \arm{P0}}} \\
Accepted, injected (placement and core cells) & 98.5 & -- & 98.0 \\
\bottomrule
\end{tabularx}
\end{table*}

The 13 real defects are those the production supervisor caught in the 28
observe-mode runs. Replaying them through the production retry (\arm{A1} at
\arm{P0}) gives a very different picture from the injected defects
(Table~\ref{tab:tier1}). The supervisor accepted the retried output for
\textbf{28.2\%} of real defects. On the same criterion, at the same placement,
it accepted \arm{A1} retries for \textbf{98.5\%} of injected defects and
\arm{A3} regenerations for 98.0\% (25.6\% on real defects). By stage, injected
reasoning defects were accepted at 100\% against 26.7\% for real ones; at the
strategy stage the figures were 88.9\% and 29.2\%. All 13 real defects fire a check class targeted by an operator; the injected report-containment class has no real-bank counterpart. Coupling at the level of the check class therefore did not
make the injected defects representative of repair.

The reason is visible in what the real defects are. The pre-study audit
(Section~\ref{sec:design-operators}) found none of the 11 audited production-caught
defects to be errors. Nine were
checker false positives with a diagnosed cause: six paraphrases that the
character-level trace matcher failed to recognise, and three single-character
identifier corruptions where the real evidence was present. The calibration of
Section~\ref{sec:results-calibration} measures the trace matcher's
false-positive rate on grounded content at 5.0\%. A retry cannot fix a defect
that is not there. When the regenerated output again paraphrases its upstream,
the same check rejects it again. On the trace class, the content-word check
that the study's scorer uses accepts 87.5\% of \arm{A1} retries, while the
production check accepts 29.2\%.

One pattern runs against the RQ1 result and should be stated rather than
smoothed. On real defects, the contentless rejection (\arm{A2}) was accepted far
more often than the production feedback (76.9\% against 28.2\%), in both
stages. A plausible reading is that naming the flagged item invites the model
to rewrite that item specifically, and it rewrites it again in words the
lexical matcher does not credit. But 13 defects in seven cells support no
inference, and the cell was pre-registered as descriptive. We report it as a
hypothesis for a larger bank of real defects. \arm{A3} returned the original
defective output in 48.7\% of trials. \arm{A1} and \arm{A2} never did.

% ---------------------------------------------------------------------
\subsection{RQ4: Repair does not become pipeline-level improvement}
\label{sec:results-rq4}

\begin{table*}[t]
\centering
\caption{Downstream propagation (90 defects with a later stage). Later stages
are re-run once each, in observe mode, from each starting point. E1: every later
stage accepted. E2: the planted \op{O5} issue appears in the final report
(20 \op{O5} defects). E3: share of the clean run's final-report claims missing
from the condition's report. ``Noise floor'' is a second clean run. Means, and
paired contrasts with 95\% intervals over defects. Usable \arm{A1}/\arm{A2} starting points comprise 70/69 reasoning and 19/19 strategy defects (89/88 total). E2 uses 20 unrepaired and 19 per retry arm.}
\label{tab:downstream}
\footnotesize
\begin{tabularx}{\linewidth}{@{}Y r r r r r@{}}
\toprule
& Clean & Noise floor & Unrepaired & \arm{A1} & \arm{A2} \\
\midrule
E1 accepted & 84.4 & -- & 80.0 & 88.8 & 88.6 \\
E2 planted issue in report & 0 & -- & \textbf{70} & 0 & 0 \\
E3 divergence from clean & -- & 2.6 & \textbf{10.8} & 6.7 & 7.4 \\
\midrule
\multicolumn{6}{@{}l}{\emph{Paired contrasts}} \\
\quad E1: unrepaired\,$-$\,clean & \multicolumn{5}{l}{$-4.4$ \ci{-14.4}{+6.7}} \\
\quad E1: \arm{A1}\,$-$\,unrepaired & \multicolumn{5}{l}{$+7.9$ \ci{-1.1}{+16.9}} \\
\quad E1: \arm{A1}\,$-$\,\arm{A2} & \multicolumn{5}{l}{$\phantom{+}0.0$ \ci{-5.7}{+5.7}} \\
\quad E3: unrepaired\,$-$\,noise floor & \multicolumn{5}{l}{$+8.2$ \ci{+5.6}{+11.0}} \\
\quad E3: \arm{A1}\,$-$\,noise floor & \multicolumn{5}{l}{$+4.1$ \ci{+1.4}{+6.8}} \\
\quad E3: \arm{A1}\,$-$\,unrepaired & \multicolumn{5}{l}{$-4.2$ \ci{-7.5}{-1.0}} \\
\quad E3: \arm{A1}\,$-$\,\arm{A2} & \multicolumn{5}{l}{$-0.6$ \ci{-3.5}{+2.3}} \\
\bottomrule
\end{tabularx}
\end{table*}

Table~\ref{tab:downstream} follows each defect into the later stages. Of the
three readings fixed before the pass, the first holds and the second holds in
part.

\paragraph{Defects propagate silently}
Pipelines carrying an \emph{unrepaired} defect were accepted by every
downstream check about as often as clean ones ($-4.4$\pp{} \ci{-14.4}{+6.7}).
Yet their final reports diverged from the clean report by 8.2\pp{} beyond the
noise floor \ci{+5.6}{+11.0}. For \op{O5}, 70\% of unrepaired runs carried the
planted, untraceable recommendation into the final report. Once a later stage
has built on a defect, the later stages' own checks do not see it.

\paragraph{Downstream outcomes after retry}
The planted issue appeared in 14/20 unrepaired reports and 0/19 reports in
each retry arm. Across all \arm{A1}-started conditions, including failed
repairs, divergence was 4.2\,pp \ci{-7.5}{-1.0} below unrepaired conditions
but 4.1\,pp \ci{+1.4}{+6.8} above the noise floor. The overall arm contrast
was 0.0\,pp on E1 and $-0.6$\,pp \ci{-3.5}{+2.3} on E3; neither establishes
a difference between arms.

\paragraph{Downstream divergence by restoration outcome}
Table~\ref{tab:downstream-split} describes post-treatment subgroups. For
reasoning defects, reports following \arm{A1} exact copies diverged 4.16\%
against their matched 2.60\% noise floor; repaired rewrites diverged 16.65\%
against their matched 3.90\% floor. The corresponding excesses for \arm{A2}
were +1.30 and +8.61\,pp. Strategy rewrites had much smaller excesses
(+0.64 and 0.00\,pp). These are observational associations, not evidence that
copying caused preservation: the groups differ after treatment and matching
noise floors does not remove selection bias. The overall arm contrast on E3
was $-0.6$\,pp \ci{-3.5}{+2.3}.

\begin{table*}[t]
\centering
\caption{Downstream report divergence (E3) partitioned by observed upstream retry outcome (RQ4).
Each upstream condition is classified as \cls{repaired\_copy}, \cls{repaired\_rewrite},
or \cls{fail} against its pre-injection clean stage output. Divergence is the share of
clean final-report claims missing; ``Noise floor'' is an in-session second clean run for
the identical defects.}
\label{tab:downstream-split}
\footnotesize
\begingroup
\setlength{\tabcolsep}{3pt}
\begin{tabular*}{\linewidth}{@{\extracolsep{\fill}}l c r rrr c rrr@{}}
\toprule
& & & \multicolumn{3}{c}{\textbf{Arm \arm{A1} (Measured Facts)}} & & \multicolumn{3}{c}{\textbf{Arm \arm{A2} (Rejection)}} \\
\cmidrule(lr){4-6} \cmidrule(lr){8-10}
Subgroup & Upstream Group & $n$ & E3 (\%) & Noise (\%) & Excess & $n$ & E3 (\%) & Noise (\%) & Excess \\
\midrule
\multirow{3}{*}{\textbf{All defects ($N=90$)}} 
  & \cls{repaired\_copy}    & 50 & 3.82 & 2.40 & \textbf{+1.43\pp} & 37 & 2.60 & 1.44 & \textbf{+1.16\pp} \\
  & \cls{repaired\_rewrite} & 31 & 9.94 & 2.27 & \textbf{+7.67\pp} & 42 & 10.20 & 4.05 & \textbf{+6.15\pp} \\
  & \cls{fail}             &  8 & 11.90 & 5.06 & \textbf{+6.85\pp} &  9 & 14.07 & 0.79 & \textbf{+13.28\pp} \\
\midrule
\multirow{3}{*}{\shortstack[l]{\textbf{Reasoning stage}\\\textbf{(\op{O1}--\op{O3}, $N=70$)}}} 
  & \cls{repaired\_copy}    & 46 (65.7\%) & 4.16 & 2.60 & \textbf{+1.55\pp} & 33 (47.8\%) & 2.91 & 1.62 & \textbf{+1.30\pp} \\
  & \cls{repaired\_rewrite} & 18 (25.7\%) & 16.65 & 3.90 & \textbf{+12.75\pp} & 30 (43.5\%) & 14.28 & 5.67 & \textbf{+8.61\pp} \\
  & \cls{fail}             &  6 (8.6\%)  & 13.10 & 6.75 & \textbf{+6.35\pp} &  6 (8.7\%)  & 15.95 & 1.19 & \textbf{+14.76\pp} \\
\midrule
\multirow{3}{*}{\shortstack[l]{\textbf{Strategy stage}\\\textbf{(\op{O5}, $N=20$)}}} 
  & \cls{repaired\_copy}    &  4 & 0.00 & 0.00 & \textbf{0.00\pp} &  4 & 0.00 & 0.00 & \textbf{0.00\pp} \\
  & \cls{repaired\_rewrite} & 13 & 0.64 & 0.00 & \textbf{+0.64\pp} & 12 & 0.00 & 0.00 & \textbf{0.00\pp} \\
  & \cls{fail}             &  2 & 8.33 & 0.00 & \textbf{+8.33\pp} &  3 & 10.32 & 0.00 & \textbf{+10.32\pp} \\
\bottomrule
\end{tabular*}
\endgroup
\end{table*}

\paragraph{The report check rewards propagation}
The 20 unrepaired \op{O5} runs split cleanly. All 14 whose report carried the
planted issue were accepted. All 6 whose report did \emph{not} carry it were
rejected by the report-containment check. That check verifies that the report
restates its upstream. When the upstream is defective, a faithful report
passes and a report that departs from the defect fails. The check measures
fidelity to upstream, not correctness.

\paragraph{Defect-level and pipeline-level outcomes are loosely coupled}
All 8 \arm{A1} runs whose defect the core had \emph{not} repaired passed every
downstream check. Of the 81 runs whose defect was repaired, 10 did not. This is
the loose coupling between failure-mode presence and task outcome that
\citet{cemri2025mast} report for multi-agent traces. Clean pipelines themselves
failed the supervisor in 15.6\% of runs (76 of 90 accepted). Across both clean repetitions, 29 of 180 conditions were rejected at the strategy traceability check.

\paragraph{Calibrated verifier acceptance}
Offline re-scoring of all 447 downstream conditions under the calibrated
content-word traceability check (\code{Jaccard}~$\ge 0.10$ replacing character-level
\code{difflib}~$\ge 0.35$ on strategy outputs) shows how acceptance changes under the alternative matcher; acceptance is not a correctness measure. Clean pipeline acceptance (E1) jumps from
\textbf{84.4\%} (76 of 90) to \textbf{96.7\%} (87 of 90): three strategy traceability flags remained under the content-word matcher; their semantic validity was not independently established.
Accepted runs under \arm{A1} rise from \textbf{88.8\%} (79 of 89) to
\textbf{97.8\%} (87 of 89), and under \arm{A2} from \textbf{88.6\%} (78 of 88)
to \textbf{98.9\%} (87 of 88). Unrepaired pipeline acceptance rises from
\textbf{80.0\%} to \textbf{93.3\%} (84 of 90; the 6 rejected runs are precisely
the 6 \op{O5} runs where the final report dropped the planted ungrounded recommendation).

% ---------------------------------------------------------------------
\subsection{RQ5: Overwrite is a cost decision, not a correctness decision}
\label{sec:results-rq5}

Routed to a retry instead of overwritten, a divergent frequency was repaired in
\textbf{69 of 69} \arm{A1} trials, for 23 of 23 defects (exact 95\% interval
over defects \ci{85.2}{100}). No trial showed collateral damage. A fresh
regeneration (\arm{A3}) repaired 97.1\% \ci{92.8}{100}, with collateral damage
in 30.4\% of trials. The overwrite is exact by construction and costs no model
call. The retry costs one call and about 22k tokens. By the reading fixed before
the freeze, the overwrite was a \emph{cost} decision, not a correctness one. The
design principle it exemplifies (any quantity the checker can compute should
not be asked of the model) rests on that zero-call argument and on exactness,
not on any failure of the retry.

% ---------------------------------------------------------------------
\subsection{Calibration of the checks}
\label{sec:results-calibration}

Two checks can be calibrated without model calls.

\paragraph{The frequency check}
The frequency check behaves exactly as specified. Sweeping \op{O4}'s divergence
over 23 bases, detection is 0\% for every divergence up to 0.04 and 100\% from
0.051. At exactly 0.05 it is 52.2\%, because the check is a strict $>0.05$ and
floating-point rounding of the injected value splits those 23 cases 12 to 11.

\paragraph{The traceability check}
The traceability check is imprecise in both directions. On the 139
recommendation issues of every clean strategy output in the 28 runs, all of
them grounded by construction, the production character-level matcher fires on
\textbf{5.0\%}: the 7 issues whose best lexical similarity to their own
upstream falls below 0.35. The content-word matcher fires on none. On 150
foreign issues placed on other businesses, the production matcher fires on
96.0\% and the content-word matcher on 100\%. The production matcher both
flags grounded paraphrase and misses foreign content. This is the measured
cause of the real-defect results in Section~\ref{sec:results-rq3}, and of the
clean-run rejections in Section~\ref{sec:results-rq4}.

A further observation concerns the pipeline, not repair. Run from its
production prompt, the report stage needed a schema-correction second call
every time: in all 447 downstream conditions, and for all 16 report-stage
defects at \arm{P0}.

% ---------------------------------------------------------------------
\subsection{RQ6: A second model}
\label{sec:results-rq6}
% Source: backend/out_study/T2L/analysis.json (repair_analysis.py --cell c7);
% gates: C7_preflight_r3/gates.json (G2 carried over from C7_preflight/);
% verification: T2L/verification.json. Exploratory throughout (A10, A11).

The second-model cell asks whether the core's answers belong to the method or
to the hosted model. Qwen2.5-7B-Instruct retried the same injected defects,
built on the hosted model's clean outputs, under the core's arms
(Section~\ref{sec:design-cells}). All of it is exploratory. G1 passed. G3
excluded four defects before any call (\op{O1} and \op{O2} on two businesses
whose reasoning prompts, at 31{,}504--31{,}916 tokens, exceeded the 28{,}608
left after the output reserve), leaving 138. G2 failed, with 8 of 10 pairs
identical, which led to A11; G4 passed, with 8 of 10 first calls valid. All
930 trials ran, 929 were usable, and every offline score agreed with its live
score.

\begin{table*}[t]
\centering
\caption{The second model (RQ6; exploratory): Qwen2.5-7B-Instruct retrying
the hosted model's injected defects at \arm{P1}, $r=3$. Rates are per-defect
means; intervals are 95\% bootstrap intervals over defects; $p$ is the exact
McNemar test on the majority dichotomy, reported descriptively.
``Identical'' compares a retried output with the clean output as canonical
JSON, the comparison behind the hosted model's identity shares.}
\label{tab:rq6}
\footnotesize
\begin{tabularx}{\linewidth}{@{}Y r r r l r@{}}
\toprule
Contrast & $N$ & \arm{A1} & \arm{A2} & Difference [95\% CI] & $p$ \\
\midrule
\textbf{All eligible defects} & \textbf{138} & \textbf{72.9} & \textbf{66.2} &
  \textbf{+6.8 \ci{+0.5}{+13.3}} & \textbf{0.052} \\
\quad first repetition only & 137 & 73.0 & 67.2 & +5.8 \ci{-0.7}{+12.4} &
  0.134 \\
\quad \op{O1} only (production-like) & 22 & 39.4 & 47.0 & $-7.6$
  \ci{-25.8}{+10.6} & 1.000 \\
\quad severe cell only (\op{O2}--\op{O7}) & 116 & 79.3 & 69.8 & +9.5
  \ci{+3.2}{+16.4} & 0.013 \\
\quad reasoning stage & 66 & 58.6 & 42.9 & & \\
\quad strategy stage & 20 & 75.0 & 80.0 & & \\
\quad report stage & 52 & 90.4 & 90.4 & & \\
\quad hosted model, same defects$^a$ & 137 & 96.8 & 95.6 & & \\
\midrule
\multicolumn{6}{@{}l}{\emph{Shares of trials (\arm{A1} 414, \arm{A2} 413) and
  of cells, \%}} \\
\quad \cls{unrepaired}: check still fires & & 7.7 & 25.9 & & \\
\quad \cls{replaced}: another passing value & & 15.7 & 6.5 & & \\
\quad identical to the clean output & & 42.0 & 39.5 & & \\
\quad collateral damage$^b$ & & 10.5 & 15.2 & & \\
\quad cells with identical raw text$^c$ & & 94.9 & 84.7 & & \\
\bottomrule
\end{tabularx}

\smallskip
\parbox{\linewidth}{\footnotesize $^a$ 137 defects for \arm{A1} and 135 for
\arm{A2}. $^b$ Of the 258 and 257 trials where it is measurable. $^c$ Of the
138 and 137 defect\,$\times$\,arm cells with three usable repetitions, those
whose raw outputs were identical in every call of all three.}
\end{table*}

\paragraph{The contrast, below the ceiling}
Both arms repaired well below the hosted model's ceiling
(Table~\ref{tab:rq6}): \textbf{72.9\%} \ci{65.5}{80.2} under measured facts and
\textbf{66.2\%} \ci{58.2}{73.9} under the rejection, a difference of
\textbf{+6.8\pp{} \ci{+0.5}{+13.3}}. The discordant defects split 16 to 6, and
the exact McNemar test gives $p=0.052$. The first repetition alone gives
+5.8\pp{} \ci{-0.7}{+12.4}, and counting the lost trial as unrepaired changes
nothing. Two of the four readings fixed in advance apply. The contrast was
measured with headroom, unlike the core's, and its interval excludes zero,
which is exploratory evidence that measured facts help this model. The
borderline test and the first-repetition interval keep that evidence a
hypothesis. The gap sits in the severe cell and at the reasoning stage, and it
is largest for fabricated identifiers (\op{O2}) and mislabelled aspects
(\op{O3}). The report stage tied, and on \op{O1}, as on the hosted model,
measured facts gained nothing.

\paragraph{How the arms fail}
After a contentless rejection the injected check still fired in \textbf{107 of
413} retries; after measured facts, in 32 of 414. Told only that its output
was wrong, the model often left the defect where it was. Told what was wrong,
it acted on the named item, though 65 times it changed it to another value that
passed the check (27 under the rejection). The hosted model left the check
firing in only 3 and 2 retries.

\paragraph{Copying, stability and regeneration}
Copy-and-correct carries over at a lower rate: 42.0\% and 39.5\% of retries
reproduced the clean output exactly, against 68\% and 53\% on the hosted model,
and still more than half of each arm's repairs did. Collateral damage was
higher, at 10.5\% and 15.2\% against 1.5\% and 0.4\%. The local runtime was
steadier than the hosted API but not bit-reproducible: across three
back-to-back repetitions, 94.9\% of \arm{A1} cells and 84.7\% of \arm{A2} cells
returned identical raw text, a stricter test than the comparison of parsed
outputs behind the hosted model's 73\%. \arm{A3}'s 38.2\% \ci{20.6}{55.9} is
not interpretable here. Every trial is scored against the hosted model's clean
output, which \arm{A1} and \arm{A2} can see and restore but a fresh draw from
another model cannot reproduce. \arm{A3} repaired no reasoning defect. Its
repairs come from the report stage, where removing an inserted item is the
repair, and all 36 of those trials needed a schema-correction call; without
them \arm{A3} repaired 4.5\%.

\paragraph{Against the hosted model}
Paired on the same defects, the second model's repair rate differed from the
hosted model's by $-24.1$\pp{} \ci{-32.4}{-16.1} under measured facts and
$-29.4$\pp{} \ci{-37.5}{-21.5} under the rejection. Family, size, quantisation and
runtime all differ, and the bases are the hosted model's. The comparison
places the ceiling in the hosted model without saying which difference
produces it.

% ---------------------------------------------------------------------
\subsection{Thinking-budget sweep (\cell{T2T})}
\label{sec:results-thinking}
% Source: backend/out_study/T2T/analysis.json (repair_analysis.py --cell thinking);
% verification: T2T/checkpoint.jsonl. Exploratory throughout (A12).

An exploratory sweep asks how configured thinking budget relates to repair and restoration.
Under Amendment~A12 (\cell{T2T}), we executed a prospective within-weights sweep
varying \texttt{reasoning\_effort} across four levels (\texttt{none}, \texttt{low},
\texttt{medium}, \texttt{high}) on \model{}, covering 45 defects ($N=45$, matching
the stratified 1/3 stride of \cell{T2P} and \cell{T2W}) with three repetitions
(1{,}224 trials, 100\% usable, 100\% offline score agreement, consuming 11.16M tokens).
Arms \arm{A1} and \arm{A2} ran at \arm{P1} on all 45 defects (540 trials per arm),
while the regeneration baseline \arm{A3} ran at \arm{P0} on 12 defects (144 trials).
Table~\ref{tab:thinking} compiles the empirical results.

\begin{table*}[t]
\centering
\caption{Within-weights thinking-budget sweep (\cell{T2T}).
Sweeping \texttt{reasoning\_effort} across four levels on \model{} ($N=45$ defects, $r=3$, 1{,}224 trials).
Rates and exact restoration $P(\text{exact copy})$ are defect-level means with 95\% cluster-bootstrap CIs.
The within-weights regeneration baseline \arm{A3}@\arm{P0} runs on $N=12$ defects (36 trials/level).
Multi-attempt counts denote trials requiring schema retry ($>1$ call).}
\label{tab:thinking}
\footnotesize
\begingroup
\setlength{\tabcolsep}{2pt}
\begin{tabular*}{\linewidth}{@{\extracolsep{\fill}}l r r c c r c c r c r r@{}}
\toprule
& Budget & Thinking tokens / & \multicolumn{3}{c}{\textbf{Binary Repair Rate (\%)}} & \multicolumn{3}{c}{\textbf{Exact Restoration $P(\text{copy})$ (\%)}} & \arm{A3}@\arm{P0} & \multicolumn{2}{c}{\textbf{Multi-Attempt}} \\
\cmidrule(lr){4-6} \cmidrule(lr){7-9} \cmidrule(lr){11-12}
Effort & Cap & Trial & \arm{A1} & \arm{A2} & Diff [95\% CI] & \arm{A1} & \arm{A2} & Diff [95\% CI] & Repair [95\% CI] & \arm{A1} & \arm{A2} \\
\midrule
\texttt{high} & 24{,}576 & 2{,}682 & 97.8 & 100.0 & $-2.2$ \ci{-6.7}{0.0} & 60.7 & 50.4 & +10.4 \ci{-3.0}{+23.0} & 75.0 \ci{50.0}{100.0} & 0 & 0 \\
\texttt{medium} & 8{,}192 & 2{,}557 & 94.8 & 94.8 & \phantom{+}0.0 \ci{-7.4}{+7.4} & 60.0 & 51.1 & +8.9 \ci{-2.2}{+20.0} & 88.9 \ci{72.2}{100.0} & 0 & 5 \\
\texttt{low} & 1{,}024 & 914 & 99.3 & 98.5 & +0.7 \ci{0.0}{+2.2} & 59.3 & 35.6 & \textbf{+23.7 \ci{+11.9}{+36.3}}$^{**}$ & 80.6 \ci{58.3}{100.0} & 0 & 4 \\
\texttt{none} & 0 & 0 & 100.0 & 100.0 & \phantom{+}0.0 \ci{0.0}{0.0} & 80.7 & 57.0 & \textbf{+23.7 \ci{+12.6}{+36.3}}$^{***}$ & \textbf{91.7 \ci{75.0}{100.0}} & 7 & 21 \\
\bottomrule
\end{tabular*}
\endgroup

\smallskip
\parbox{\linewidth}{\footnotesize $p$-values from exact McNemar tests on discordant defect majorities: all binary repair contrasts have $p=1.000$; for exact restoration, $p=0.302$ (\texttt{high}), $p=0.754$ (\texttt{medium}), $^{**}p=0.007$ (\texttt{low}), and $^{***}p=0.001$ (\texttt{none}).}
\end{table*}

At \texttt{none}, all 135 trials in each feedback arm were scored repaired.
This demonstrates the observed performance with a zero thinking-token budget,
not a single forward pass: seven \arm{A1} and 21 \arm{A2} trials required
schema-correction calls, compared with none at \texttt{high}. Mean \arm{A1}
latency was 7.17 seconds at \texttt{none} versus 17.44 seconds at \texttt{high}.
The exact-restoration contrast varied nonmonotonically across levels, reaching
+23.7\,pp \ci{+12.6}{+36.3} at \texttt{none}; the exploratory effort-step
slope was +5.48\,pp \ci{+0.44}{+10.59}, whereas the log-budget slope interval
included zero. \arm{A3} repair also varied nonmonotonically (75.0\%, 88.9\%,
80.6\%, 91.7\% from high to none). These results do not identify internal
reasoning mechanisms.
